\section{Costate identities for reverse switches}\label{app:reverse-estimates}

This appendix gives the explicit formulas behind \cref{sec:periods,sec:reverse}: the one-piece
costate map, the two reverse-four identities, the sign of the cost derivative $N$, and the
elementary inequalities of the three-arc lemma, the switch norms and the Casimir cap. Each
identity below is an identity of rational functions in the displayed variables and the
logarithms that occur; it is verified by clearing denominators and expanding.

\subsection{The one-piece costate map}\label{app:rev-one-arc}

On a piece with active margin $a$, growing margin $b$ and decreasing margin $c$ at its start
(\cref{lem:one-arc}), the inhomogeneous terms of \cref{lem:costate-transport}(b) at the factor
$q$ are
\begin{align*}
  \mathrm{Act}(a,b,c,q)&=\tfrac{\sqrt3a}{2}\bigl((a+c)(1-q^{-1})-(a+b)(q-1)\bigr),\\
  \mathrm{AG}(a,b,q)&=\tfrac{\sqrt3a}{2}\bigl(a(q^2-1)-(a+b)(q^2-q)\bigr),\\
  \mathrm{AD}(a,c,q)&=\tfrac{\sqrt3a}{2}\bigl((a+c)(q^{-1}-q^{-2})-a(1-q^{-2})\bigr),\\
  \mathrm{BG}(a,b,q)&=\sqrt3a\bigl((a+b)\,q\log q-2a(q-1)\bigr),\\
  \mathrm{BD}(a,c,q)&=\sqrt3a\,q^{-1}\bigl(2a(q-1)-(a+c)\log q\bigr).
\end{align*}
The logarithms come from the cost term $-\frac32[J,X]$ of the costate
equation, integrated against the exponentially moving margins.

These formulas give the switching functions of the two inactive labels inside the piece. Let
$i$ be the active label, and at the start of the piece let $Y=\alpha_i+\alpha_{i+1}$,
$Y'=\alpha_i+\alpha_{i-1}$, $B=\beta_i+\beta_{i+1}$, $B'=\beta_i+\beta_{i-1}$. Since $\chi_i=0$ on
the piece, at the intermediate factor $s\in[1,q]$
\[
  \chi_{i+1}(s)=2s^2Y+sB+\sqrt3\,a\,s\,F(s),\qquad
  \chi_{i-1}(s)=\frac{2Y'}{s^2}+\frac{B'}{s}-\frac{\sqrt3\,a}{s}\,h(s),
\]
where
\begin{align*}
  F(s)&=(a+b)\log s-(a+b)(s-1)+a\,\frac{(s-1)^2}{s},\\
  h(s)&=(a+c)\log s-a(s-1)-c\,(1-s^{-1}),
\end{align*}
because $2\mathrm{AG}+\mathrm{BG}=\sqrt3asF$ and $2\mathrm{AD}+\mathrm{BD}=-\sqrt3a\,h/s$. Both $F$ and
$h$ vanish to second order at $s=1$: $F(1)=F'(1)=h(1)=h'(1)=0$. When the growing label was active
on the previous piece, $\chi_{i+1}(1)=2Y+B=0$ and
\[
  \chi_{i+1}(s)=2s\Bigl(Y(s-1)+\frac{\sqrt3a}{2}F(s)\Bigr);
\]
when the decreasing label was, $B'=-2Y'$ and
\[
  \chi_{i-1}(s)=-\frac2s\Bigl(Y'(1-s^{-1})+\frac{\sqrt3a}{2}h(s)\Bigr).
\]
These are the forms used in \cref{lem:switch-norms}(a) and in \cref{lem:three-arc}. In each case the derivative at $s=1$ is $\pm2Y$ or $\pm2Y'$, which is how
the sign condition at the start of a piece fixes the sign of $Y$ or $Y'$.

\subsection{The reverse-four identities}\label{app:rev-identities}

In the notation of \cref{sec:rev-four}, let $x_j,y_j$ be the coordinates $\alpha,\beta$ at the
switch $\tau_{k+1}$ in the label order $(i,i+1,i-1)$, so that $j=0,1,2$ stands for the labels
$i,i+1,i-1$, and let $u_j,v_j$ and $r_j,s_j$ be those at $\tau_{k+2}$ and $\tau_{k+3}$ in the same
order. The transport over piece $k+1$ (active margin $Z$, growing margin $P$, decreasing margin
$Q$, factor $b$) and over piece $k+2$ (active margin $A$, growing margin $Z$, decreasing margin
$B$, factor $c$) gives
\begin{align*}
  u_2&=x_2+\mathrm{Act}(Z,P,Q,b),\\
  u_0&=b^2(x_2+x_0)+\mathrm{AG}(Z,P,b)-u_2,\\
  u_1&=b^{-2}(x_2+x_1)+\mathrm{AD}(Z,Q,b)-u_2,\\
  r_1&=u_1+\mathrm{Act}(A,Z,B,c),\\
  r_0&=c^{-2}(u_1+u_0)+\mathrm{AD}(A,B,c)-r_1,\\
  r_2&=c^2(u_1+u_2)+\mathrm{AG}(A,Z,c)-r_1,
\end{align*}
and the same with $(\alpha,\mathrm{Act},\mathrm{AG},\mathrm{AD})$ replaced by
$(\beta,-2\mathrm{Act},\mathrm{BG},\mathrm{BD})$ and the powers $b^{\pm2},c^{\pm2}$ by
$b^{\pm1},c^{\pm1}$. The five relations are
\begin{gather*}
  h_1=2x_0+y_0,\quad h_2=2x_2+y_2,\quad h_3=2u_1+v_1,\quad h_4=2r_0+s_0,\\
  h_5=y_0(Q+Z)+y_1(P+Z)+y_2(P+Q),
\end{gather*}
all equal to $0$: the active switching functions at the three switches, and
$\ip{\Lambda_R}{X}=0$ at $\tau_{k+1}$. Then, with $D=A(c-1)+Z(b-1)$ and as identities in all the
variables,
\begin{align*}
  B(c-1)&(2x_1+y_1)+bD(2u_0+v_0)+\sqrt3ZN(A,B,Z,b,c)\\
  &=b^2(c-1)(Q-Zb+Z)\,h_1+(b^2-1)(c-1)(Pb+Q)\,h_2\\
  &\quad+b(c-1)(Pb^2+Zb^2+Zbc-Zc-Z)\,h_3+Zbc^2(b-1)\,h_4\\
  &\quad-b(b-1)(c-1)\,h_5,
\end{align*}
and
\begin{align*}
  (b-1)&B(2r_2+s_2)+cD(2u_0+v_0)+\sqrt3AN(Z,B,A,c,b)\\
  &=bc(c-1)(Q-Zb+Z)\,h_1\\
  &\quad+\tfrac{c(b-1)}{b}\bigl(Pb^2c+Qbc-Qb+Qc-Q+Zb^2-Zbc+Zc-Z\bigr)\,h_2\\
  &\quad+(b-1)(c^2-1)(Pb+Zb+Zc-Z)\,h_3+Zc^3(b-1)\,h_4\\
  &\quad-c(b-1)(c-1)\,h_5,
\end{align*}
with $A=(Z+Q)/b-Z$ and $B=(Z+P)b-Z$. Setting the five
relations to zero gives \cref{lem:rev-identities}.

\subsection{The sign of the cost derivative}\label{app:rev-N}

Let $\eta(t)=t-1-\log t$ and $\zeta(t)=t\log t-t+1$. For $t>1$ both are positive.

\begin{lemma}[Pad\'e bound]\label{lem:rev-pade}
For $t>1$, $3(t^2-1)<(t^2+4t+1)\log t$; equivalently $(2t+1)\,\eta(t)<(t+2)\,\zeta(t)$. In
particular $\eta(t)<\zeta(t)$.
\end{lemma}

\begin{proof}
The function $\psi(t)=\log t-3(t^2-1)/(t^2+4t+1)$ vanishes at $1$ and
$\psi'(t)=(t-1)^4/\bigl(t(t^2+4t+1)^2\bigr)>0$ for $t>1$. The second form follows from
$(t+2)\zeta(t)-(2t+1)\eta(t)=(t^2+4t+1)\log t-3(t^2-1)$, and the third because $(2t+1)/(t+2)>1$
and $\eta(t)>0$.
\end{proof}

\begin{proof}[Proof of \cref{lem:rev-N}]
Recall~\eqref{eq:rev-N} and the hypotheses $b,c>1$, $A,Z>0$, $A\le Z$, $bA+(b-1)Z\le B$ and
$(2c-1)A\le B$. In particular $B>cA>0$.

The first bracket of $N$ is negative: $Ac\,\eta(c)<Ac\,\zeta(c)<B\,\zeta(c)$.

The second bracket, $T(B)=bA(B+bZ)\,\eta(b)-B(bA+Z)\,\zeta(b)$, is affine in $B$ with slope
$bA\,\eta(b)-(bA+Z)\,\zeta(b)<0$, so $T(B)\le T(B_0)$ with $B_0=bA+(b-1)Z$. At $B_0$ we have
$B_0+bZ=bA+(2b-1)Z$ and $B_0(bA+Z)=b^2A^2+b^2AZ+(b-1)Z^2>0$, and \cref{lem:rev-pade} at $t=b$ gives
\begin{align*}
  (b+2)\,T(B_0)&<\eta(b)\bigl[(b+2)\,bA\bigl(bA+(2b-1)Z\bigr)\\
  &\qquad\qquad-(2b+1)\bigl(b^2A^2+b^2AZ+(b-1)Z^2\bigr)\bigr]\\
  &=-(b-1)\,\eta(b)\bigl(b^2A^2+2bZ(Z-A)+Z^2\bigr),
\end{align*}
which is $\le0$ because $Z\ge A$. So $T(B)<0$, and
$N=b(b-1)A\cdot(\text{first bracket})+(c-1)\,T(B)<0$.
\end{proof}

\subsection{Elementary inequalities}\label{app:rev-elementary}

\runin{Logarithms.} For $q>1$, $2(q-1)<(q+1)\log q$ and $q-1<q\log q$: both differences vanish
at $q=1$, with derivatives
$\log q+q^{-1}-1>0$ and $\log q>0$. They give the signs in the proof of \cref{lem:three-arc}.

\runin{The turn secant.} With $f(x)=x-x^{-1}-\log x$,
\[
  f'(x)-\tfrac34=\frac{(x-2)^2}{4x^2},\qquad 1-f'(x)=\frac{x-1}{x^2}.
\]
So every secant slope of $f$ on $(1,\infty)$ lies in $[3/4,1)$, and the switch norm at a turn,
$d=\frac18\times$(secant slope), satisfies $3/32\le d<1/8$. At a forward switch the lower bound
$d_k\ge3/32$ comes instead from
$j_k/2-3/32=(a_k-a_{k+1})^2/(32s_k^2)\ge0$.

\runin{The cross pairing bound.} With $a,p,n>0$ and $j(a,b)=(a^2+ab+b^2)/(4(a+b)^2)$,
\begin{multline*}
  \frac{1}{2a}\Bigl((a+p)\,\frac{j(p,a)}{2}+(a+n)\,\frac{j(a,n)}{2}
  -\frac{G_a(a,p)+G_a(a,n)}{16}\Bigr)\\
  =\frac{1}{16}\Bigl(-1+\frac{2a}{a+p}+\frac{2a}{a+n}-\frac{a^2}{a^2+4}
  +\log\frac{(a+p)(a+n)}{a^2+4}\Bigr).
\end{multline*}
This is the quantity $X_k$ of \cref{lem:rev-cross-signs}: by
\cref{lem:stationarity}, at a same-sign piece $e_k$ is an increasing affine function of
$d_{k-1}$ and $d_k$, and $X_k$ is its value at $d=j/2$.

\runin{The Casimir arithmetic.} The proof of \cref{cor:casimir-cap} uses $\Aoct<3.127$ and $\pi>3.14$. From $4\pi\kappa\le(2\kappa+3)\cost$ and
$\cost<3.127$,
\[
  \kappa<\frac{3\cdot3.127}{4\cdot3.14-2\cdot3.127}=1.48763\ldots<1.4877,
\]
and $1.4877^2/12=0.18443\ldots<97969/529200=0.18512\ldots$.
